% ============================================================
%  封面 / Cover page
%  设计:深海军蓝底 + 金色细规 + 数据曲线motif
%  在 main.tex 中以 % ============================================================
%  封面 / Cover page
%  设计:深海军蓝底 + 金色细规 + 数据曲线motif
%  在 main.tex 中以 % ============================================================
%  封面 / Cover page
%  设计:深海军蓝底 + 金色细规 + 数据曲线motif
%  在 main.tex 中以 % ============================================================
%  封面 / Cover page
%  设计:深海军蓝底 + 金色细规 + 数据曲线motif
%  在 main.tex 中以 \input{cover} 引入
% ============================================================

% 配色在 main.tex 的 preamble 中统一定义(coverbg / gold / cyanx / ink / muted / faint)

\begin{titlepage}
\thispagestyle{empty}
\begin{tikzpicture}[remember picture, overlay, shift={(current page.south west)}, x=1cm, y=1cm]

  %% ---------- 背景 ----------
  \fill[coverbg] (0,0) rectangle (21,29.7);

  %% 极淡点阵网格(科技感底纹)
  \foreach \x in {0.6,1.2,...,20.4}{%
    \foreach \y in {0.6,1.2,...,29.4}{%
      \fill[white, opacity=0.05] (\x,\y) circle (0.017);}}

  %% 右上角装饰:同心弧
  \foreach \r in {3.2,4.0,4.8,5.6}{%
    \draw[cyanx, opacity=0.10, line width=0.4pt] (21,29.7) circle (\r);}

  %% ---------- 页眉 ----------
  \node[anchor=north west, text=muted, font=\sffamily\fontsize{9}{11}\selectfont]
    at (2.4,28.4) {AI 金融研学实训营 \; · \; AI 篇 \; · \; 主题四};
  \node[anchor=north east, text=gold, font=\ttfamily\fontsize{9}{11}\selectfont]
    at (18.6,28.4) {GROUP 14};
  \draw[gold, line width=0.6pt] (2.4,27.85) -- (18.6,27.85);

  %% ---------- 主标题 ----------
  \node[anchor=north west, text=ink, align=left,
        font=\heiti\fontsize{34}{46}\selectfont]
    at (2.4,25.4) {可解释 Alpha 因子\\生成智能体};

  %% 金色短粗强调条
  \fill[gold] (2.4,19.55) rectangle (4.5,19.68);

  %% ---------- 副标题 ----------
  \node[anchor=north west, text=ink, align=left,
        font=\sffamily\fontsize{13}{20}\selectfont]
    at (2.4,19.15) {点时间对齐、受限表达式引擎\\与确定性回测的可信底座};

  \node[anchor=north west, text=muted, align=left,
        font=\sffamily\itshape\fontsize{9.5}{13.5}\selectfont]
    at (2.4,16.1)
    {An Explainable Alpha-Factor Generation Agent:\\
     Point-in-Time Alignment, a Restricted Expression Engine,\\
     and a Fully Deterministic Backtest};

  %% ---------- 四步闭环 chips ----------
  \begin{scope}[shift={(2.4,13.2)}]
    \foreach \i/\txt in {0/生成, 1/校验, 3/反馈}{%
      \node[anchor=west, text=muted, font=\sffamily\fontsize{10}{10}\selectfont]
        at ({\i*3.5},0) {\txt};}
    % 回测:高亮为金色 + 外框,强调"纯代码"
    \node[anchor=west, text=gold, font=\sffamily\bfseries\fontsize{10}{10}\selectfont]
      at (7.0,0) {回测};
    \draw[gold, opacity=0.55, line width=0.5pt, rounded corners=2pt]
      (6.75,-0.32) rectangle (8.35,0.34);
    \node[anchor=north, text=gold, opacity=0.9,
          font=\ttfamily\fontsize{7}{7}\selectfont] at (7.55,-0.46) {no LLM};
    % 箭头连接(避开回测外框)
    \draw[-{Latex[length=1.5mm]}, color=faint, line width=0.5pt] (0.95,0) -- (2.85,0);
    \draw[-{Latex[length=1.5mm]}, color=faint, line width=0.5pt] (4.45,0) -- (6.55,0);
    \draw[-{Latex[length=1.5mm]}, color=faint, line width=0.5pt] (8.55,0) -- (10.35,0);
    % 回环:圆角回到"生成"
    \draw[-{Latex[length=1.5mm]}, color=faint, line width=0.5pt, rounded corners=4pt]
      (11.3,-0.28) -- (11.3,-1.15) -- (0.28,-1.15) -- (0.28,-0.28);
    \node[anchor=north, text=muted, opacity=0.6,
          font=\sffamily\fontsize{7.5}{7.5}\selectfont] at (5.8,-1.18) {迭代闭环};
  \end{scope}

  %% ---------- 数据曲线 motif ----------
  \begin{scope}
    % 基线与刻度
    \draw[faint, opacity=0.55, line width=0.4pt] (2.4,8.0) -- (18.6,8.0);
    \foreach \x in {2.4,4.7,...,18.6}{%
      \draw[faint, opacity=0.35, line width=0.3pt] (\x,8.0) -- (\x,10.9);}

    % 面积填充
    \fill[cyanx, opacity=0.07]
      (2.4,8.0) -- (2.4,8.35) -- (4.0,8.75) -- (5.6,8.5) -- (7.2,9.25)
      -- (8.8,9.65) -- (10.4,9.3) -- (12.0,10.0) -- (13.6,9.7)
      -- (15.2,10.35) -- (16.8,10.75) -- (18.6,11.05) -- (18.6,8.0) -- cycle;

    % 曲线
    \draw[cyanx, line width=1.1pt, line join=round]
      (2.4,8.35) -- (4.0,8.75) -- (5.6,8.5) -- (7.2,9.25)
      -- (8.8,9.65) -- (10.4,9.3) -- (12.0,10.0) -- (13.6,9.7)
      -- (15.2,10.35) -- (16.8,10.75) -- (18.6,11.05);

    % 节点
    \foreach \p in {(2.4,8.35),(4.0,8.75),(5.6,8.5),(7.2,9.25),(8.8,9.65),%
                    (10.4,9.3),(12.0,10.0),(13.6,9.7),(15.2,10.35),%
                    (16.8,10.75),(18.6,11.05)}{%
      \fill[coverbg] \p circle (0.075);
      \draw[cyanx, line width=0.7pt] \p circle (0.075);}

    \node[anchor=north west, text=muted, opacity=0.75,
          font=\ttfamily\fontsize{7.5}{7.5}\selectfont]
      at (2.4,7.85) {cumulative rank IC · out-of-sample};
  \end{scope}

  %% ---------- 研究团队 ----------
  \draw[faint, line width=0.5pt] (2.4,6.35) -- (18.6,6.35);
  \node[anchor=north west, text=gold,
        font=\ttfamily\fontsize{8}{8}\selectfont] at (2.4,6.1) {RESEARCH TEAM};
  \node[anchor=north west, text=ink, align=left,
        font=\sffamily\fontsize{11.5}{17}\selectfont] at (2.4,5.5)
    {周梓煜\;\textcolor{muted}{\fontsize{9}{9}\selectfont（组长）}\quad
     张心湉\quad 许艺馨\quad 张修博\quad 刘迅羽\quad 王佳祺};

  %% ---------- 底栏 ----------
  \fill[coverbg2] (0,0) rectangle (21,2.6);
  \draw[gold, opacity=0.7, line width=0.6pt] (0,2.6) -- (21,2.6);

  \node[anchor=west, text=ink, font=\sffamily\fontsize{10}{10}\selectfont]
    at (2.4,1.55) {LLM 提出假设，确定性代码作出裁决。};
  \node[anchor=west, text=muted, font=\ttfamily\fontsize{8}{8}\selectfont]
    at (2.4,0.92) {The LLM proposes hypotheses; deterministic code renders the verdict.};

  \node[anchor=east, text=muted, font=\ttfamily\fontsize{8.5}{11}\selectfont]
    at (18.6,1.55) {\today};
  \node[anchor=east, text=muted, opacity=0.7,
        font=\ttfamily\fontsize{8.5}{11}\selectfont]
    at (18.6,0.95) {v1.0 · synthetic-data build};

\end{tikzpicture}
\end{titlepage}
 引入
% ============================================================

% 配色在 main.tex 的 preamble 中统一定义(coverbg / gold / cyanx / ink / muted / faint)

\begin{titlepage}
\thispagestyle{empty}
\begin{tikzpicture}[remember picture, overlay, shift={(current page.south west)}, x=1cm, y=1cm]

  %% ---------- 背景 ----------
  \fill[coverbg] (0,0) rectangle (21,29.7);

  %% 极淡点阵网格(科技感底纹)
  \foreach \x in {0.6,1.2,...,20.4}{%
    \foreach \y in {0.6,1.2,...,29.4}{%
      \fill[white, opacity=0.05] (\x,\y) circle (0.017);}}

  %% 右上角装饰:同心弧
  \foreach \r in {3.2,4.0,4.8,5.6}{%
    \draw[cyanx, opacity=0.10, line width=0.4pt] (21,29.7) circle (\r);}

  %% ---------- 页眉 ----------
  \node[anchor=north west, text=muted, font=\sffamily\fontsize{9}{11}\selectfont]
    at (2.4,28.4) {AI 金融研学实训营 \; · \; AI 篇 \; · \; 主题四};
  \node[anchor=north east, text=gold, font=\ttfamily\fontsize{9}{11}\selectfont]
    at (18.6,28.4) {GROUP 14};
  \draw[gold, line width=0.6pt] (2.4,27.85) -- (18.6,27.85);

  %% ---------- 主标题 ----------
  \node[anchor=north west, text=ink, align=left,
        font=\heiti\fontsize{34}{46}\selectfont]
    at (2.4,25.4) {可解释 Alpha 因子\\生成智能体};

  %% 金色短粗强调条
  \fill[gold] (2.4,19.55) rectangle (4.5,19.68);

  %% ---------- 副标题 ----------
  \node[anchor=north west, text=ink, align=left,
        font=\sffamily\fontsize{13}{20}\selectfont]
    at (2.4,19.15) {点时间对齐、受限表达式引擎\\与确定性回测的可信底座};

  \node[anchor=north west, text=muted, align=left,
        font=\sffamily\itshape\fontsize{9.5}{13.5}\selectfont]
    at (2.4,16.1)
    {An Explainable Alpha-Factor Generation Agent:\\
     Point-in-Time Alignment, a Restricted Expression Engine,\\
     and a Fully Deterministic Backtest};

  %% ---------- 四步闭环 chips ----------
  \begin{scope}[shift={(2.4,13.2)}]
    \foreach \i/\txt in {0/生成, 1/校验, 3/反馈}{%
      \node[anchor=west, text=muted, font=\sffamily\fontsize{10}{10}\selectfont]
        at ({\i*3.5},0) {\txt};}
    % 回测:高亮为金色 + 外框,强调"纯代码"
    \node[anchor=west, text=gold, font=\sffamily\bfseries\fontsize{10}{10}\selectfont]
      at (7.0,0) {回测};
    \draw[gold, opacity=0.55, line width=0.5pt, rounded corners=2pt]
      (6.75,-0.32) rectangle (8.35,0.34);
    \node[anchor=north, text=gold, opacity=0.9,
          font=\ttfamily\fontsize{7}{7}\selectfont] at (7.55,-0.46) {no LLM};
    % 箭头连接(避开回测外框)
    \draw[-{Latex[length=1.5mm]}, color=faint, line width=0.5pt] (0.95,0) -- (2.85,0);
    \draw[-{Latex[length=1.5mm]}, color=faint, line width=0.5pt] (4.45,0) -- (6.55,0);
    \draw[-{Latex[length=1.5mm]}, color=faint, line width=0.5pt] (8.55,0) -- (10.35,0);
    % 回环:圆角回到"生成"
    \draw[-{Latex[length=1.5mm]}, color=faint, line width=0.5pt, rounded corners=4pt]
      (11.3,-0.28) -- (11.3,-1.15) -- (0.28,-1.15) -- (0.28,-0.28);
    \node[anchor=north, text=muted, opacity=0.6,
          font=\sffamily\fontsize{7.5}{7.5}\selectfont] at (5.8,-1.18) {迭代闭环};
  \end{scope}

  %% ---------- 数据曲线 motif ----------
  \begin{scope}
    % 基线与刻度
    \draw[faint, opacity=0.55, line width=0.4pt] (2.4,8.0) -- (18.6,8.0);
    \foreach \x in {2.4,4.7,...,18.6}{%
      \draw[faint, opacity=0.35, line width=0.3pt] (\x,8.0) -- (\x,10.9);}

    % 面积填充
    \fill[cyanx, opacity=0.07]
      (2.4,8.0) -- (2.4,8.35) -- (4.0,8.75) -- (5.6,8.5) -- (7.2,9.25)
      -- (8.8,9.65) -- (10.4,9.3) -- (12.0,10.0) -- (13.6,9.7)
      -- (15.2,10.35) -- (16.8,10.75) -- (18.6,11.05) -- (18.6,8.0) -- cycle;

    % 曲线
    \draw[cyanx, line width=1.1pt, line join=round]
      (2.4,8.35) -- (4.0,8.75) -- (5.6,8.5) -- (7.2,9.25)
      -- (8.8,9.65) -- (10.4,9.3) -- (12.0,10.0) -- (13.6,9.7)
      -- (15.2,10.35) -- (16.8,10.75) -- (18.6,11.05);

    % 节点
    \foreach \p in {(2.4,8.35),(4.0,8.75),(5.6,8.5),(7.2,9.25),(8.8,9.65),%
                    (10.4,9.3),(12.0,10.0),(13.6,9.7),(15.2,10.35),%
                    (16.8,10.75),(18.6,11.05)}{%
      \fill[coverbg] \p circle (0.075);
      \draw[cyanx, line width=0.7pt] \p circle (0.075);}

    \node[anchor=north west, text=muted, opacity=0.75,
          font=\ttfamily\fontsize{7.5}{7.5}\selectfont]
      at (2.4,7.85) {cumulative rank IC · out-of-sample};
  \end{scope}

  %% ---------- 研究团队 ----------
  \draw[faint, line width=0.5pt] (2.4,6.35) -- (18.6,6.35);
  \node[anchor=north west, text=gold,
        font=\ttfamily\fontsize{8}{8}\selectfont] at (2.4,6.1) {RESEARCH TEAM};
  \node[anchor=north west, text=ink, align=left,
        font=\sffamily\fontsize{11.5}{17}\selectfont] at (2.4,5.5)
    {周梓煜\;\textcolor{muted}{\fontsize{9}{9}\selectfont（组长）}\quad
     张心湉\quad 许艺馨\quad 张修博\quad 刘迅羽\quad 王佳祺};

  %% ---------- 底栏 ----------
  \fill[coverbg2] (0,0) rectangle (21,2.6);
  \draw[gold, opacity=0.7, line width=0.6pt] (0,2.6) -- (21,2.6);

  \node[anchor=west, text=ink, font=\sffamily\fontsize{10}{10}\selectfont]
    at (2.4,1.55) {LLM 提出假设，确定性代码作出裁决。};
  \node[anchor=west, text=muted, font=\ttfamily\fontsize{8}{8}\selectfont]
    at (2.4,0.92) {The LLM proposes hypotheses; deterministic code renders the verdict.};

  \node[anchor=east, text=muted, font=\ttfamily\fontsize{8.5}{11}\selectfont]
    at (18.6,1.55) {\today};
  \node[anchor=east, text=muted, opacity=0.7,
        font=\ttfamily\fontsize{8.5}{11}\selectfont]
    at (18.6,0.95) {v1.0 · synthetic-data build};

\end{tikzpicture}
\end{titlepage}
 引入
% ============================================================

% 配色在 main.tex 的 preamble 中统一定义(coverbg / gold / cyanx / ink / muted / faint)

\begin{titlepage}
\thispagestyle{empty}
\begin{tikzpicture}[remember picture, overlay, shift={(current page.south west)}, x=1cm, y=1cm]

  %% ---------- 背景 ----------
  \fill[coverbg] (0,0) rectangle (21,29.7);

  %% 极淡点阵网格(科技感底纹)
  \foreach \x in {0.6,1.2,...,20.4}{%
    \foreach \y in {0.6,1.2,...,29.4}{%
      \fill[white, opacity=0.05] (\x,\y) circle (0.017);}}

  %% 右上角装饰:同心弧
  \foreach \r in {3.2,4.0,4.8,5.6}{%
    \draw[cyanx, opacity=0.10, line width=0.4pt] (21,29.7) circle (\r);}

  %% ---------- 页眉 ----------
  \node[anchor=north west, text=muted, font=\sffamily\fontsize{9}{11}\selectfont]
    at (2.4,28.4) {AI 金融研学实训营 \; · \; AI 篇 \; · \; 主题四};
  \node[anchor=north east, text=gold, font=\ttfamily\fontsize{9}{11}\selectfont]
    at (18.6,28.4) {GROUP 14};
  \draw[gold, line width=0.6pt] (2.4,27.85) -- (18.6,27.85);

  %% ---------- 主标题 ----------
  \node[anchor=north west, text=ink, align=left,
        font=\heiti\fontsize{34}{46}\selectfont]
    at (2.4,25.4) {可解释 Alpha 因子\\生成智能体};

  %% 金色短粗强调条
  \fill[gold] (2.4,19.55) rectangle (4.5,19.68);

  %% ---------- 副标题 ----------
  \node[anchor=north west, text=ink, align=left,
        font=\sffamily\fontsize{13}{20}\selectfont]
    at (2.4,19.15) {点时间对齐、受限表达式引擎\\与确定性回测的可信底座};

  \node[anchor=north west, text=muted, align=left,
        font=\sffamily\itshape\fontsize{9.5}{13.5}\selectfont]
    at (2.4,16.1)
    {An Explainable Alpha-Factor Generation Agent:\\
     Point-in-Time Alignment, a Restricted Expression Engine,\\
     and a Fully Deterministic Backtest};

  %% ---------- 四步闭环 chips ----------
  \begin{scope}[shift={(2.4,13.2)}]
    \foreach \i/\txt in {0/生成, 1/校验, 3/反馈}{%
      \node[anchor=west, text=muted, font=\sffamily\fontsize{10}{10}\selectfont]
        at ({\i*3.5},0) {\txt};}
    % 回测:高亮为金色 + 外框,强调"纯代码"
    \node[anchor=west, text=gold, font=\sffamily\bfseries\fontsize{10}{10}\selectfont]
      at (7.0,0) {回测};
    \draw[gold, opacity=0.55, line width=0.5pt, rounded corners=2pt]
      (6.75,-0.32) rectangle (8.35,0.34);
    \node[anchor=north, text=gold, opacity=0.9,
          font=\ttfamily\fontsize{7}{7}\selectfont] at (7.55,-0.46) {no LLM};
    % 箭头连接(避开回测外框)
    \draw[-{Latex[length=1.5mm]}, color=faint, line width=0.5pt] (0.95,0) -- (2.85,0);
    \draw[-{Latex[length=1.5mm]}, color=faint, line width=0.5pt] (4.45,0) -- (6.55,0);
    \draw[-{Latex[length=1.5mm]}, color=faint, line width=0.5pt] (8.55,0) -- (10.35,0);
    % 回环:圆角回到"生成"
    \draw[-{Latex[length=1.5mm]}, color=faint, line width=0.5pt, rounded corners=4pt]
      (11.3,-0.28) -- (11.3,-1.15) -- (0.28,-1.15) -- (0.28,-0.28);
    \node[anchor=north, text=muted, opacity=0.6,
          font=\sffamily\fontsize{7.5}{7.5}\selectfont] at (5.8,-1.18) {迭代闭环};
  \end{scope}

  %% ---------- 数据曲线 motif ----------
  \begin{scope}
    % 基线与刻度
    \draw[faint, opacity=0.55, line width=0.4pt] (2.4,8.0) -- (18.6,8.0);
    \foreach \x in {2.4,4.7,...,18.6}{%
      \draw[faint, opacity=0.35, line width=0.3pt] (\x,8.0) -- (\x,10.9);}

    % 面积填充
    \fill[cyanx, opacity=0.07]
      (2.4,8.0) -- (2.4,8.35) -- (4.0,8.75) -- (5.6,8.5) -- (7.2,9.25)
      -- (8.8,9.65) -- (10.4,9.3) -- (12.0,10.0) -- (13.6,9.7)
      -- (15.2,10.35) -- (16.8,10.75) -- (18.6,11.05) -- (18.6,8.0) -- cycle;

    % 曲线
    \draw[cyanx, line width=1.1pt, line join=round]
      (2.4,8.35) -- (4.0,8.75) -- (5.6,8.5) -- (7.2,9.25)
      -- (8.8,9.65) -- (10.4,9.3) -- (12.0,10.0) -- (13.6,9.7)
      -- (15.2,10.35) -- (16.8,10.75) -- (18.6,11.05);

    % 节点
    \foreach \p in {(2.4,8.35),(4.0,8.75),(5.6,8.5),(7.2,9.25),(8.8,9.65),%
                    (10.4,9.3),(12.0,10.0),(13.6,9.7),(15.2,10.35),%
                    (16.8,10.75),(18.6,11.05)}{%
      \fill[coverbg] \p circle (0.075);
      \draw[cyanx, line width=0.7pt] \p circle (0.075);}

    \node[anchor=north west, text=muted, opacity=0.75,
          font=\ttfamily\fontsize{7.5}{7.5}\selectfont]
      at (2.4,7.85) {cumulative rank IC · out-of-sample};
  \end{scope}

  %% ---------- 研究团队 ----------
  \draw[faint, line width=0.5pt] (2.4,6.35) -- (18.6,6.35);
  \node[anchor=north west, text=gold,
        font=\ttfamily\fontsize{8}{8}\selectfont] at (2.4,6.1) {RESEARCH TEAM};
  \node[anchor=north west, text=ink, align=left,
        font=\sffamily\fontsize{11.5}{17}\selectfont] at (2.4,5.5)
    {周梓煜\;\textcolor{muted}{\fontsize{9}{9}\selectfont（组长）}\quad
     张心湉\quad 许艺馨\quad 张修博\quad 刘迅羽\quad 王佳祺};

  %% ---------- 底栏 ----------
  \fill[coverbg2] (0,0) rectangle (21,2.6);
  \draw[gold, opacity=0.7, line width=0.6pt] (0,2.6) -- (21,2.6);

  \node[anchor=west, text=ink, font=\sffamily\fontsize{10}{10}\selectfont]
    at (2.4,1.55) {LLM 提出假设，确定性代码作出裁决。};
  \node[anchor=west, text=muted, font=\ttfamily\fontsize{8}{8}\selectfont]
    at (2.4,0.92) {The LLM proposes hypotheses; deterministic code renders the verdict.};

  \node[anchor=east, text=muted, font=\ttfamily\fontsize{8.5}{11}\selectfont]
    at (18.6,1.55) {\today};
  \node[anchor=east, text=muted, opacity=0.7,
        font=\ttfamily\fontsize{8.5}{11}\selectfont]
    at (18.6,0.95) {v1.0 · synthetic-data build};

\end{tikzpicture}
\end{titlepage}
 引入
% ============================================================

% 配色在 main.tex 的 preamble 中统一定义(coverbg / gold / cyanx / ink / muted / faint)

\begin{titlepage}
\thispagestyle{empty}
\begin{tikzpicture}[remember picture, overlay, shift={(current page.south west)}, x=1cm, y=1cm]

  %% ---------- 背景 ----------
  \fill[coverbg] (0,0) rectangle (21,29.7);

  %% 极淡点阵网格(科技感底纹)
  \foreach \x in {0.6,1.2,...,20.4}{%
    \foreach \y in {0.6,1.2,...,29.4}{%
      \fill[white, opacity=0.05] (\x,\y) circle (0.017);}}

  %% 右上角装饰:同心弧
  \foreach \r in {3.2,4.0,4.8,5.6}{%
    \draw[cyanx, opacity=0.10, line width=0.4pt] (21,29.7) circle (\r);}

  %% ---------- 页眉 ----------
  \node[anchor=north west, text=muted, font=\sffamily\fontsize{9}{11}\selectfont]
    at (2.4,28.4) {AI 金融研学实训营 \; · \; AI 篇 \; · \; 主题四};
  \node[anchor=north east, text=gold, font=\ttfamily\fontsize{9}{11}\selectfont]
    at (18.6,28.4) {GROUP 14};
  \draw[gold, line width=0.6pt] (2.4,27.85) -- (18.6,27.85);

  %% ---------- 主标题 ----------
  \node[anchor=north west, text=ink, align=left,
        font=\heiti\fontsize{34}{46}\selectfont]
    at (2.4,25.4) {可解释 Alpha 因子\\生成智能体};

  %% 金色短粗强调条
  \fill[gold] (2.4,19.55) rectangle (4.5,19.68);

  %% ---------- 副标题 ----------
  \node[anchor=north west, text=ink, align=left,
        font=\sffamily\fontsize{13}{20}\selectfont]
    at (2.4,19.15) {点时间对齐、受限表达式引擎\\与确定性回测的可信底座};

  \node[anchor=north west, text=muted, align=left,
        font=\sffamily\itshape\fontsize{9.5}{13.5}\selectfont]
    at (2.4,16.1)
    {An Explainable Alpha-Factor Generation Agent:\\
     Point-in-Time Alignment, a Restricted Expression Engine,\\
     and a Fully Deterministic Backtest};

  %% ---------- 四步闭环 chips ----------
  \begin{scope}[shift={(2.4,13.2)}]
    \foreach \i/\txt in {0/生成, 1/校验, 3/反馈}{%
      \node[anchor=west, text=muted, font=\sffamily\fontsize{10}{10}\selectfont]
        at ({\i*3.5},0) {\txt};}
    % 回测:高亮为金色 + 外框,强调"纯代码"
    \node[anchor=west, text=gold, font=\sffamily\bfseries\fontsize{10}{10}\selectfont]
      at (7.0,0) {回测};
    \draw[gold, opacity=0.55, line width=0.5pt, rounded corners=2pt]
      (6.75,-0.32) rectangle (8.35,0.34);
    \node[anchor=north, text=gold, opacity=0.9,
          font=\ttfamily\fontsize{7}{7}\selectfont] at (7.55,-0.46) {no LLM};
    % 箭头连接(避开回测外框)
    \draw[-{Latex[length=1.5mm]}, color=faint, line width=0.5pt] (0.95,0) -- (2.85,0);
    \draw[-{Latex[length=1.5mm]}, color=faint, line width=0.5pt] (4.45,0) -- (6.55,0);
    \draw[-{Latex[length=1.5mm]}, color=faint, line width=0.5pt] (8.55,0) -- (10.35,0);
    % 回环:圆角回到"生成"
    \draw[-{Latex[length=1.5mm]}, color=faint, line width=0.5pt, rounded corners=4pt]
      (11.3,-0.28) -- (11.3,-1.15) -- (0.28,-1.15) -- (0.28,-0.28);
    \node[anchor=north, text=muted, opacity=0.6,
          font=\sffamily\fontsize{7.5}{7.5}\selectfont] at (5.8,-1.18) {迭代闭环};
  \end{scope}

  %% ---------- 数据曲线 motif ----------
  \begin{scope}
    % 基线与刻度
    \draw[faint, opacity=0.55, line width=0.4pt] (2.4,8.0) -- (18.6,8.0);
    \foreach \x in {2.4,4.7,...,18.6}{%
      \draw[faint, opacity=0.35, line width=0.3pt] (\x,8.0) -- (\x,10.9);}

    % 面积填充
    \fill[cyanx, opacity=0.07]
      (2.4,8.0) -- (2.4,8.35) -- (4.0,8.75) -- (5.6,8.5) -- (7.2,9.25)
      -- (8.8,9.65) -- (10.4,9.3) -- (12.0,10.0) -- (13.6,9.7)
      -- (15.2,10.35) -- (16.8,10.75) -- (18.6,11.05) -- (18.6,8.0) -- cycle;

    % 曲线
    \draw[cyanx, line width=1.1pt, line join=round]
      (2.4,8.35) -- (4.0,8.75) -- (5.6,8.5) -- (7.2,9.25)
      -- (8.8,9.65) -- (10.4,9.3) -- (12.0,10.0) -- (13.6,9.7)
      -- (15.2,10.35) -- (16.8,10.75) -- (18.6,11.05);

    % 节点
    \foreach \p in {(2.4,8.35),(4.0,8.75),(5.6,8.5),(7.2,9.25),(8.8,9.65),%
                    (10.4,9.3),(12.0,10.0),(13.6,9.7),(15.2,10.35),%
                    (16.8,10.75),(18.6,11.05)}{%
      \fill[coverbg] \p circle (0.075);
      \draw[cyanx, line width=0.7pt] \p circle (0.075);}

    \node[anchor=north west, text=muted, opacity=0.75,
          font=\ttfamily\fontsize{7.5}{7.5}\selectfont]
      at (2.4,7.85) {cumulative rank IC · out-of-sample};
  \end{scope}

  %% ---------- 研究团队 ----------
  \draw[faint, line width=0.5pt] (2.4,6.35) -- (18.6,6.35);
  \node[anchor=north west, text=gold,
        font=\ttfamily\fontsize{8}{8}\selectfont] at (2.4,6.1) {RESEARCH TEAM};
  \node[anchor=north west, text=ink, align=left,
        font=\sffamily\fontsize{11.5}{17}\selectfont] at (2.4,5.5)
    {周梓煜\;\textcolor{muted}{\fontsize{9}{9}\selectfont（组长）}\quad
     张心湉\quad 许艺馨\quad 张修博\quad 刘迅羽\quad 王佳祺};

  %% ---------- 底栏 ----------
  \fill[coverbg2] (0,0) rectangle (21,2.6);
  \draw[gold, opacity=0.7, line width=0.6pt] (0,2.6) -- (21,2.6);

  \node[anchor=west, text=ink, font=\sffamily\fontsize{10}{10}\selectfont]
    at (2.4,1.55) {LLM 提出假设，确定性代码作出裁决。};
  \node[anchor=west, text=muted, font=\ttfamily\fontsize{8}{8}\selectfont]
    at (2.4,0.92) {The LLM proposes hypotheses; deterministic code renders the verdict.};

  \node[anchor=east, text=muted, font=\ttfamily\fontsize{8.5}{11}\selectfont]
    at (18.6,1.55) {\today};
  \node[anchor=east, text=muted, opacity=0.7,
        font=\ttfamily\fontsize{8.5}{11}\selectfont]
    at (18.6,0.95) {v1.0 · synthetic-data build};

\end{tikzpicture}
\end{titlepage}
